% Page 090

\seite{90}

\margmark{20. Februar  S}%
eit Mitte Dezember hat der Winter eine\ob
furchtbare Schärfe angenommen. Nur wenige\ob
Tage mit normaler Temperatur sind zu verzeich=\ob
nen. Die Morgenfrühen eine ununterbrechbaren\ob
Kälteherrschaft. Einig 69\unsicher und tiefer! 20 und\ob
mehr Grad Kälte sind keine Seltenheit mehr.\ob
Fast alle Wasserleistungen, Hausleitungen, sind\ob
eingefroren. Küche und Flüsse bilden nur einzel=\ob
ne Eisfläche. Das Wild verhungert in Feld und Wald.\ob
Die ältesten Dorfbewohner haben einen solchen\ob
langen strengen Winter noch nicht erlebt.

In der Generalversammlung des Kriegervereins\ob
wurde Kameradschaft der Pfarrei Bettingen unter=\ob
dessen Ringo-Sefferweich (1.2.29) wieder ein=\ob
stimmig zum Vorsitzenden gewählt.

\margmark{10. März     E}%
ndlich scheint die Eiskälte gebrochen\ob
zu sein. Nur langsam weichen Rhein und Eis.\ob
Die Ost-Lage sind noch vollständig vereist.\ob
Da die Abdeckeisfolge bis im Mai gefroren ist,\ob
stehen die Wasserleitungsrohre immer noch\ob
unter Eis, und das Wasser muß 3 Tage schon\ob
monatelang aus Cisternen, Dorfbrunnen\ob
bei die Angela geholt werden.

Nachdem die Masernzeitdemie kaum\ob
vorüber war, setzte eine Grippeepidemie ein,\ob
von der fast sämtliche Dorfinsassen in Mit=\ob
leidenschaft gezogen wurden, da sind 55% der\ob
Schulkinder so krank waren, mußte auf Ver=\ob
anlassung des Kreisarztes, d.\ h. Ermäßen, die\ob
Schule vom 6.-10.\ März geschlossen bleiben.\ob
Sind 2 Mitglieder der Kriegsgemeinde, welche\ob
Opfer der Grippe, und im außerdem mußte der\ob
Bäcker zwei seines Raumsteuern mit militäri=\ob
schen Ehren zu Grabe geleiten (Alter 73 u.\ 56 Jahre), und\ob
ein 83 Jahre Goßmuth\unsicher – beide von Seffern\unsicher.
\pend
